% New section. Place after Section~\ref{sec:q4} and before
% "Verification of the implementation".
% Requires the generated table gen_budget_tradeoff.tex.

\section{Budget and the block-dimension trade-off}
\label{sec:budget}

\subsection{Motivation}

Section~\ref{sec:classfunction} showed that at $M=\dim\Qcal$ every admissible
pool spans the whole quotient, so the sector partition cannot distinguish
selection rules. Two questions follow. First, whether the agreement between
cost-ranked rules reported in Section~\ref{sec:q4} survives when that
degeneracy is removed. Second, what is paid, in block dimension, for reducing
the budget. Both are addressed by repeating the campaign of
Section~\ref{sec:q4} at $M=\dim\Qcal-1$, that is $M=3$ for $\mathrm{H_2O}$ and $M=5$
for $\mathrm{N_2}$, with all other settings unchanged.

\subsection{The agreement is not an artefact of the budget}

At $M=\dim\Qcal-1$ the spanned subspace is a proper subspace of $\Qcal$, so
different rules \emph{may} select different subspaces. The exhaustive and
iterative rules nonetheless span the identical subspace at all $42$ points of
the reduced-budget campaign, and again return identical $K$. The agreement of
Section~\ref{sec:q4} is therefore a property of greedy optimality on the
matroid, as Proposition~\ref{prop:coset} and
Section~\ref{sec:classfunction}(ii) predict, and not a consequence of the
degenerate budget.

We record this explicitly because the converse inference is available and
wrong: observing agreement at $M=\dim\Qcal$ alone would not have distinguished
the structural explanation from the degeneracy, since at that budget the
partition is rule-independent by construction and the span test cannot fail.

\subsection{Cost of reducing the budget}

\begin{table}[htbp]
\centering
\begin{tabular}{llrrrr}
\hline
molecule & budget & $M$ & $\Dmax$ & retained dim & $K$ range\\
\hline
$\mathrm{H_2O}$ & $n-\rsp$ & 4 & 21 & 133 & 4--8\\
$\mathrm{H_2O}$ & $n-\rsp-1$ & 3 & 27,31 & 133 & 1--3\\
$\mathrm{N_2}$ & $n-\rsp$ & 6 & 120 & 1824 & 4--26\\
$\mathrm{N_2}$ & $n-\rsp-1$ & 5 & 156 & 1824 & 3--25\\
\hline
\end{tabular}

\caption{Effect of reducing the budget by one generator, at the exhaustive rule
with the NC cost. $\Dmax$ is the largest joint exact-and-approximate sector;
``retained dim'' is the total dimension of the sectors surviving the exact
filter; the $K$ range is over the geometries of the scan. Where two values of
$\Dmax$ appear, they occur at different geometries.}
\label{tab:tradeoff}
\end{table}

Three features of Table~\ref{tab:tradeoff} are relevant.

\paragraph{The retained dimension is invariant.}
The dimension of the retained exact sector, $\sum_{s}\dim(s)$ over sectors
surviving the exact filter, is unchanged between budgets. This is required:
that dimension is fixed by $\Ecal$ and the chosen exact sector alone, and the
LAS rows only subdivide it. Its invariance is used here as a consistency check
on the two campaigns rather than as a result.

\paragraph{Removing one generator coarsens the partition.}
\begin{sloppypar}
$\Dmax$, the largest joint exact-and-approximate sector and hence the largest
block that must be diagonalised, rises from $120$ to $156$ for $\mathrm{N_2}$, a
factor of $1.300$ at every geometry, and from $21$ to $27$ for $\mathrm{H_2O}$ at nine
of ten geometries, a factor of $1.286$. The exception is $\mathrm{H_2O}$ at
$R=2.5$\,\AA, where $\Dmax=31$ (factor $1.476$); this is also the only $\mathrm{H_2O}$
geometry at which the reduced budget does not reach the floor discussed below,
and the two observations have the same origin --- the leading sector there is
not large enough to absorb the target state. This increase is the direct cost
of the reduced budget.
\end{sloppypar}

\paragraph{The coupled dimension falls, but not monotonically.}
For $\mathrm{N_2}$ the coupled dimension decreases at eight of eleven geometries, the
range falling from $4$--$26$ to $3$--$25$, since a coarser partition places more
of the target state in the leading sector. It increases at the remaining three,
$R=1.7$--$1.9$\,\AA\ ($23\to25$, $18\to19$, $14\to15$). Because the partition at
fixed budget is identical across rules, an increase in $K$ at reduced budget
cannot originate in the partition; the only remaining channel is the orbital
optimisation converging to a different stationary point.

We do not attribute it further. In particular these three geometries are
\emph{not} those at which the quota rule departs most from the cost-ranked
rules: the largest $|\Delta K|$ between rules occurs at $R=1.2$\,\AA\ at
$M=n-\rsp$ and at $R=1.7$\,\AA\ at $M=n-\rsp-1$, and at $R=1.8$\,\AA\ the two
rules agree exactly at the reduced budget. The budget-induced and rule-induced
departures therefore do not localise to the same geometries, and there is no
evidence here that they are the same phenomenon. Both are candidates for the
multi-start study of Section~\ref{sec:conclusions}, which is the only way to
establish whether either exceeds the scatter of the optimiser.

\paragraph{$\mathrm{H_2O}$ reaches the floor.}
At $M=3$ the decoupled energy is already within $\varepsilon_E$ of
$E_{\mathrm{target}}$ at nine of ten geometries, so $K=1$ and no coupling is
required. No comparative information about selection rules is available from
these points, and they are excluded from the Q4(ii) reading for that reason.

\subsection{A consequence for budget selection}

For $\mathrm{H_2O}$ in this basis the reduced budget is not merely tolerable but
preferable on both axes that the procedure exposes. Dropping from $M=4$ to
$M=3$ removes one generator, and in exchange $K$ falls from the range $4$--$8$
to $1$ at nine of ten geometries and $3$ at the tenth, while $\Dmax$ rises from
$21$ to $27$ (to $31$ at $R=2.5$\,\AA). The coupled problem is thereby
eliminated at nine of ten geometries at a cost of one moderately larger block.

Whether that exchange is favourable depends on the cost model, and the present
data do not settle it: the coupled stage scales with $K$ and the number of
retained sectors, the decoupled stage with $\Dmax$ per block, and the two are
not reduced to a common measure anywhere in this report. What the data do
establish is that the trade exists and is not one-sided --- reducing the budget
is not simply a degradation --- and that the budget is therefore a parameter
worth scanning rather than fixing at the bound $M=n-\rsp$. For $\mathrm{N_2}$
the same reduction lowers $K$ at eight of eleven geometries and raises it at
three, so no comparable statement holds there.
